\SourceLeaf{336}%
\textsc{Bessel} in Toisen angegebenen Werthe von \(a\), Art. 5), wenn \(r\) ursprünglich in Toisen ausgedrückt war, zu dem Logarithmen hinzusetzen haben
\[
0{,}2897827662
\]
oder, wenn \(r\) ursprünglich in französischen gesetzlichen Metern gegeben war, wird von dem Logarithmen subtrahirt werden müssen
\[
0{,}0000371638
\]

Die Glieder, welche die Factoren (3), (4) u.\,s.\,w. enthalten, können als Correctionen betrachtet werden, durch welche die genäherten Logarithmen \(\log\tau,\log\beta,\log\lambda\) in die berichtigten \(\log t,\log b,\log l\) verwandelt werden. Diese Correctionen sind in allen Fällen, für welche unsre Methode angewandt werden soll, nur sehr kleine Decimalbrüche, und da jene Logarithmen in der Regel siebenzifrig gerechnet werden, so ist es bequem, auch jene Correctionen sofort in Einheiten der siebenten Decimale ausgedrückt zu erhalten. Dies geschieht, indem man den Coëfficienten (3), (4) u.\,s.\,w. anstatt der im vorhergehenden Art. angegebenen Werthe zehnmillionenmal grössere beilegt, oder ihre Logarithmen um sieben Einheiten vergrössert. Auf diese Weise sind sie in unsrer Tafel angesetzt, und so wird denn auch
\[
\begin{aligned}
\log(7)&=4{,}62872\ (-10)\\
\log3(7)&=5{,}10584\ (-10)
\end{aligned}
\]
gesetzt werden. Übrigens sind auch so noch (3), (4), (5), (6), eben so wie (1) und (2) ächte Brüche, oder ihre Logarithmen an sich negativ: in der Tafel stehen sie aber nach üblicher Art, indem sämmtlichen Logarithmen 10 Einheiten geborgt sind.

\Article{35}
Von der Benutzung unsrer Formeln zur Auflösung der zu Anfang dieser Abhandlung aufgestellten Aufgabe gilt nun alles, was oben (Art. 20) in Beziehung auf dieselbe Aufgabe für die Kugelfläche gesagt ist, fast unverändert und unter geringen Modificationen. Bezeichnet man die wirklich gegebenen Grössen, nemlich die Breite und das Azimuth an dem ersten Orte mit \(B^0\) und \(T^0\), so wird man zuerst, von einem genäherten Werthe von \(T\) ausgehend (wofür man in Er%
